\section{Discussion}
\label{sec:discussion}

\subsection{Why fidelity falls, and when that is the right trade}
\label{sec:psnr_tradeoff}
Peak signal to noise ratio usually falls slightly when the wrapper is applied, and the size of the
drop differs by backbone. This is not an accident of tuning. It follows from what the two objectives
reward. A denoiser trained on squared error is at a local optimum of that error, so any change to its
output moves away from that optimum. The wrapper makes changes that raise contrast and boundary
definition, and those changes are not free in squared error terms. The question is therefore not how
to avoid the drop but how to decide when it is worth accepting.

Two things decide it. A change of a few tenths of a decibel is smaller than the variation between two
acquisitions of the same eye, so it does not change what a downstream algorithm can do. And the
quantities that improve, namely boundary definition and tissue to background separation, are what a
grader actually uses to place a segmentation boundary or judge a fluid pocket.

The backbone that loses most fidelity is the one that had suppressed the most contrast to begin with.
A transformer with windowed attention smooths aggressively inside a window, which is efficient for
squared error and costly for local contrast. The wrapper restores that contrast, and restoring it
necessarily moves the pixels further from the smoothed optimum. The pattern is therefore informative
rather than embarrassing. The size of the correction tracks how much the backbone had given away.

For a deployment that must not lose fidelity at all, the safety layer already provides the control.
Tightening the third constraint of Section~\ref{sec:safety} reduces the admitted change, and the
result moves smoothly toward the backbone output. What the method does not offer is a free lunch in
which every measure improves at once, and we prefer to say so.

\subsection{What the formal results buy, and what they do not}
Theorem~\ref{thm:stability} is exact and it is about the code that ran. The constant follows from the
trained weights and Section~\ref{sec:theory_check} measures the realised ratio against it. That is a
different claim from a bound that holds vacuously because a clip enforces it.

Theorem~\ref{thm:cnr} is a sufficient condition rather than a guarantee of improvement. It says which
three quantities have to behave for contrast not to fall, and the method is built to make all three
behave. Reporting how often they actually hold, rather than asserting that they must, is the honest
form of that claim.

\subsection{Why the correction layer is worth its complexity}
A plain network with the same parameter budget can also improve clinical measures, as
Section~\ref{sec:matched} shows. What it cannot do is say why it changed a pixel. Every change made
here can be traced to a named property, a location where that property was measured to fail, a rule
that fired, and a safety decision that admitted the result. That trace is the product, and the
comparison at matched complexity is what shows the trace is not bought by giving up performance.

\subsection{Limitations}
The correction is bounded by construction, so a backbone that has destroyed a structure cannot be
made to recover it. The formal results are narrow in the way stated at the end of
Section~\ref{sec:stability}. The confidence head is the one part that must be refitted for a new
backbone, and Section~\ref{sec:transfer_cost} reports what that costs rather than treating it as
free. The clinical measures used here are established proxies, not outcomes, and no clinician read
the images, so the safety analysis of Section~\ref{sec:safety_check} is quantitative evidence rather
than clinical validation. Finally, the external datasets are small, sixteen and eighteen images, so
the transfer results should be read as evidence of behaviour across scanners rather than as a
precise estimate of the gain on any one device.
